% fig_architecture_overview.tex
% Figure 1 – PP-VAE-Hformer overall architecture (Hformer Fig-1 style)
% % fig_architecture_overview.tex
% Figure 1 – PP-VAE-Hformer overall architecture (Hformer Fig-1 style)
% % fig_architecture_overview.tex
% Figure 1 – PP-VAE-Hformer overall architecture (Hformer Fig-1 style)
% % fig_architecture_overview.tex
% Figure 1 – PP-VAE-Hformer overall architecture (Hformer Fig-1 style)
% \input{fig_architecture_overview} from main.tex  (no \documentclass)
\begin{landscape}
\begin{figure}[p]
\centering
\resizebox{0.98\linewidth}{!}{%
\begin{tikzpicture}[
  >=Stealth,
  every node/.style={font=\small},
  % ── block styles ───────────────────────────────────────────────────
  imgblk/.style={inner sep=1pt, outer sep=0pt},
  stemblk/.style={draw=gray!70, fill=cyan!8, rounded corners=3pt,
                   minimum width=2.0cm, minimum height=1.2cm,
                   align=center, line width=0.7pt},
  encblk/.style={draw=blue!55!black, fill=blue!9, rounded corners=3pt,
                  minimum width=2.3cm, minimum height=1.3cm,
                  align=center, line width=0.7pt},
  botblk/.style={draw=teal!70!black, fill=teal!12, rounded corners=3pt,
                  minimum width=2.6cm, minimum height=1.3cm,
                  align=center, line width=1.0pt},
  vaeblk/.style={draw=purple!60!black, fill=purple!7, rounded corners=3pt,
                  minimum width=3.4cm, minimum height=1.6cm,
                  align=center, line width=1.0pt, dashed},
  decblk/.style={draw=orange!65!black, fill=orange!9, rounded corners=3pt,
                  minimum width=2.3cm, minimum height=1.3cm,
                  align=center, line width=0.7pt},
  headblk/.style={draw=teal!55!black, fill=teal!8, rounded corners=3pt,
                   minimum width=1.9cm, minimum height=0.95cm,
                   align=center, line width=0.7pt},
  mcblk/.style={draw=purple!50!black, fill=purple!6, rounded corners=3pt,
                 minimum width=2.4cm, minimum height=0.9cm,
                 align=center, line width=0.8pt, dashed},
  dimtag/.style={font=\fontsize{7}{8}\selectfont, text=gray!65, align=center},
  % ── arrow styles ───────────────────────────────────────────────────
  arrmain/.style={->, line width=1.2pt, gray!55},
  arrskip/.style={->, line width=1.0pt, blue!55, dashed},
  arrbyp/.style={->, line width=0.8pt, gray!40, dashed},
  arrvae/.style={->, line width=0.9pt, purple!55},
]

% ────────────────────────────────────────────────────────────────────
% INPUT THUMBNAIL
% ────────────────────────────────────────────────────────────────────
\node[imgblk] (inp) at (0,8.8)
  {\includegraphics[width=1.6cm,height=1.6cm,keepaspectratio]%
    {figures/arch/arch_input_noisy.png}};
\node[dimtag,below=2pt of inp] {Noisy CXR\\[-1pt]$1{\times}256{\times}256$};

% 1-ch feature grid at input
\begin{scope}[shift={(1.1,8.5)}]
  \fill[gray!25] (0.04,0.04) rectangle (0.84,0.84);
  \fill[gray!35,draw=gray!60,line width=0.4pt] (0,0) rectangle (0.8,0.8);
  \draw[gray!50,very thin,step=0.2] (0,0) grid (0.8,0.8);
\end{scope}

% ────────────────────────────────────────────────────────────────────
% STEM
% ────────────────────────────────────────────────────────────────────
\node[stemblk] (stem) at (2.7,8.8)
  {\textbf{Stem}\\{\fontsize{7}{8}\selectfont Conv $3{\times}3$, pad\,1\\$1\!\to\!64\,\mathrm{ch}$}};

% 64-ch feature map after stem (256²)
\begin{scope}[shift={(3.85,8.5)}]
  \fill[blue!15] (0.04,0.04) rectangle (0.84,0.84);
  \fill[blue!20,draw=blue!45,line width=0.4pt] (0,0) rectangle (0.8,0.8);
  \draw[blue!35,very thin,step=0.2] (0,0) grid (0.8,0.8);
  \node[dimtag,above=0pt] at (0.4,0.85) {$64{\times}256^2$};
\end{scope}

% ────────────────────────────────────────────────────────────────────
% ENCODER
% ────────────────────────────────────────────────────────────────────
\node[encblk] (enc1) at (5.2,8.8)
  {\textbf{E1}\\{\fontsize{7}{8}\selectfont HB${\times}2+$Down}};
\node[dimtag,above=2pt of enc1] {$64{\times}256^2$};

% 128-ch feature map after E1 + downsample (128²)
\begin{scope}[shift={(6.45,8.5)}]
  \fill[blue!18] (0.04,0.04) rectangle (0.69,0.69);
  \fill[blue!24,draw=blue!50,line width=0.4pt] (0,0) rectangle (0.65,0.65);
  \draw[blue!40,very thin,step=0.1625] (0,0) grid (0.65,0.65);
  \node[dimtag,above=0pt] at (0.325,0.7) {$128{\times}128^2$};
\end{scope}

\node[encblk] (enc2) at (7.4,6.3)
  {\textbf{E2}\\{\fontsize{7}{8}\selectfont HB${\times}2+$Down}};
\node[dimtag,above left=1pt and 0pt of enc2] {$128{\times}128^2$};

% 256-ch feature map after E2 + downsample (64²)
\begin{scope}[shift={(8.65,6.05)}]
  \fill[blue!22] (0.04,0.04) rectangle (0.54,0.54);
  \fill[blue!28,draw=blue!55,line width=0.4pt] (0,0) rectangle (0.5,0.5);
  \draw[blue!45,very thin,step=0.125] (0,0) grid (0.5,0.5);
  \node[dimtag,above=0pt] at (0.25,0.55) {$256{\times}64^2$};
\end{scope}

\node[encblk] (enc3) at (9.6,3.8)
  {\textbf{E3}\\{\fontsize{7}{8}\selectfont HB${\times}2+$Down}};
\node[dimtag,above left=1pt and 0pt of enc3] {$256{\times}64^2$};

% 512-ch feature map after E3 + downsample (32²)
\begin{scope}[shift={(10.85,3.6)}]
  \fill[teal!20] (0.04,0.04) rectangle (0.42,0.42);
  \fill[teal!26,draw=teal!55,line width=0.4pt] (0,0) rectangle (0.38,0.38);
  \draw[teal!45,very thin,step=0.095] (0,0) grid (0.38,0.38);
  \node[dimtag,above=0pt] at (0.19,0.43) {$512{\times}32^2$};
\end{scope}

% ────────────────────────────────────────────────────────────────────
% BOTTLENECK
% ────────────────────────────────────────────────────────────────────
\node[botblk] (bot) at (11.9,1.3)
  {\textbf{Bottleneck}\\{\fontsize{7}{8}\selectfont HB${\times}2$}};
\node[dimtag,below=2pt of bot] {$512{\times}32^2$};

% 512-ch feature map after bottleneck
\begin{scope}[shift={(13.3,1.1)}]
  \fill[teal!20] (0.04,0.04) rectangle (0.42,0.42);
  \fill[teal!26,draw=teal!55,line width=0.4pt] (0,0) rectangle (0.38,0.38);
  \draw[teal!45,very thin,step=0.095] (0,0) grid (0.38,0.38);
\end{scope}

% ────────────────────────────────────────────────────────────────────
% VAE BOTTLENECK
% ────────────────────────────────────────────────────────────────────
\node[vaeblk] (vae) at (15.6,1.3)
  {\textbf{VAE Bottleneck}\\
   {\fontsize{7}{8}\selectfont
    $\mathbf{z}_\mu,\,\mathbf{z}_{\log\sigma^2}$\;(Conv $1{\times}1$)\\[2pt]
    $\mathbf{z}{=}\mathbf{z}_\mu{+}\boldsymbol{\varepsilon}{\odot}
     \exp\!\bigl(\tfrac{1}{2}\mathbf{z}_{\log\sigma^2}\bigr)$\\
    $\boldsymbol{\varepsilon}{\sim}\mathcal{N}(\mathbf{0},\mathbf{I})$}};

% ────────────────────────────────────────────────────────────────────
% DECODER
% ────────────────────────────────────────────────────────────────────
\node[decblk] (dec3) at (18.8,3.8)
  {\textbf{D3}\\{\fontsize{7}{8}\selectfont Up$+s_3+$HB${\times}2$}};
\node[dimtag,above right=1pt and 0pt of dec3] {$256{\times}64^2$};

% 256-ch feature map after D3
\begin{scope}[shift={(20.05,3.6)}]
  \fill[orange!20] (0.04,0.04) rectangle (0.54,0.54);
  \fill[orange!26,draw=orange!55,line width=0.4pt] (0,0) rectangle (0.5,0.5);
  \draw[orange!45,very thin,step=0.125] (0,0) grid (0.5,0.5);
  \node[dimtag,above=0pt] at (0.25,0.55) {$256{\times}64^2$};
\end{scope}

\node[decblk] (dec2) at (21.0,6.3)
  {\textbf{D2}\\{\fontsize{7}{8}\selectfont Up$+s_2+$HB${\times}2$}};
\node[dimtag,above right=1pt and 0pt of dec2] {$128{\times}128^2$};

% 128-ch feature map after D2
\begin{scope}[shift={(22.25,6.05)}]
  \fill[orange!18] (0.04,0.04) rectangle (0.69,0.69);
  \fill[orange!23,draw=orange!50,line width=0.4pt] (0,0) rectangle (0.65,0.65);
  \draw[orange!40,very thin,step=0.1625] (0,0) grid (0.65,0.65);
  \node[dimtag,above=0pt] at (0.325,0.7) {$128{\times}128^2$};
\end{scope}

\node[decblk] (dec1) at (23.2,8.8)
  {\textbf{D1}\\{\fontsize{7}{8}\selectfont Up$+s_1+$HB${\times}2$}};
\node[dimtag,above=2pt of dec1] {$64{\times}256^2$};

% ────────────────────────────────────────────────────────────────────
% DUAL HEAD + OUTPUT THUMBNAILS
% ────────────────────────────────────────────────────────────────────
\node[headblk] (hmu) at (25.2,9.5)
  {$\hat{\mu}_y$\\{\fontsize{7}{8}\selectfont Conv $3{\times}3$}};
\node[headblk,draw=orange!60!black,fill=orange!8] (hsig) at (25.2,8.1)
  {$\log\hat{\sigma}^2_a$\\{\fontsize{7}{8}\selectfont Conv $3{\times}3$}};
\node[dimtag,above=2pt of hmu] {Dual Head};

\node[imgblk] (outmu) at (27.0,9.5)
  {\includegraphics[width=1.4cm,height=1.4cm,keepaspectratio]%
    {figures/arch/arch_output_mu.png}};
\node[dimtag,below=1pt of outmu] {$\hat{\mu}_y$};

\node[imgblk] (outsig) at (27.0,8.1)
  {\includegraphics[width=1.4cm,height=1.4cm,keepaspectratio]%
    {figures/arch/arch_output_aleat.png}};
\node[dimtag,below=1pt of outsig] {$\hat{\sigma}^2_a$};

% ────────────────────────────────────────────────────────────────────
% K=20 EPISTEMIC BRANCH
% ────────────────────────────────────────────────────────────────────
\node[mcblk] (mc) at (15.6,-0.3)
  {{\fontsize{7}{8}\selectfont $K{=}20$ stochastic passes}};
\node[imgblk] (outepi) at (18.8,-0.3)
  {\includegraphics[width=1.4cm,height=1.4cm,keepaspectratio]%
    {figures/arch/arch_output_epist.png}};
\node[dimtag,below=1pt of outepi] {$\hat{\sigma}^2_e$\,(epistemic)};

% ────────────────────────────────────────────────────────────────────
% ARROWS – main flow
% ────────────────────────────────────────────────────────────────────
\draw[arrmain] (inp.east)          -- (stem.west);
\draw[arrmain] (stem.east)         -- (enc1.west);
\draw[arrmain] (enc1.south east)   -- (enc2.north west);
\draw[arrmain] (enc2.south east)   -- (enc3.north west);
\draw[arrmain] (enc3.south east)   -- (bot.north west);
\draw[arrmain] (bot.east)          -- (vae.west);
\draw[arrmain] (vae.east)          -| (dec3.south);
\draw[arrmain] (dec3.north east)   -- (dec2.south west);
\draw[arrmain] (dec2.north east)   -- (dec1.south west);
\draw[arrmain] (dec1.east) -- ++(0.4,0) |- (hmu.west);
\draw[arrmain] (dec1.east) -- ++(0.4,0) |- (hsig.west);
\draw[arrmain] (hmu.east)          -- (outmu.west);
\draw[arrmain] (hsig.east)         -- (outsig.west);

% Arms A-D deterministic bypass (skips VAE)
\draw[arrbyp] (bot.north) -- ++(0,1.5) -|
  node[above,pos=0.3,font=\fontsize{6}{7}\selectfont,text=gray!55]
    {Arms\,A--D\;(no\,VAE)}
  (dec3.west);

% VAE → epistemic MC branch
\draw[arrvae] (vae.south) -- (mc.north);
\draw[arrvae] (mc.east)   -- (outepi.west);
\node[dimtag,right=2pt of outepi] {Var of $K$ predictions};

% ────────────────────────────────────────────────────────────────────
% SKIP CONNECTIONS (dashed blue, routed above encoder blocks)
% enc1.x=5.2, dec1.x=23.2  → dx=18.0, route at y=10.05
% enc2.x=7.4, dec2.x=21.0  → dx=13.6, route at y=7.4
% enc3.x=9.6, dec3.x=18.8  → dx=9.2,  route at y=4.85
% ────────────────────────────────────────────────────────────────────
\draw[arrskip]
  (enc1.north) -- ++(0,0.6) --
  node[above,font=\fontsize{7}{8}\selectfont,black!70]{$s_1$}
  ++(18.0,0) -- (dec1.north);

\draw[arrskip]
  (enc2.north) -- ++(0,0.45) --
  node[above,font=\fontsize{7}{8}\selectfont,black!70]{$s_2$}
  ++(13.6,0) -- (dec2.north);

\draw[arrskip]
  (enc3.north) -- ++(0,0.4) --
  node[above,font=\fontsize{7}{8}\selectfont,black!70]{$s_3$}
  ++(9.2,0) -- (dec3.north);

\end{tikzpicture}%
}
\caption[PP-VAE-Hformer overall architecture]{%
  \textbf{PP-VAE-Hformer encoder--decoder architecture.}
  A noisy paediatric CXR ($1{\times}256{\times}256$) enters a stem convolution
  and three encoder stages (\textbf{E1--E3}), each comprising two HybridBlocks
  (HB${\times}2$) followed by strided downsampling; channel depth expands
  $64\!\to\!128\!\to\!256\,\text{ch}$ while spatial resolution halves.
  The Bottleneck applies two further HybridBlocks at $512{\times}32^2$.
  The probabilistic VAE Bottleneck (dashed border) maps features to a spatial
  latent via the reparameterisation trick; Arms\,A--D bypass this stage
  (grey dashed arrow).
  Three symmetric decoder stages (\textbf{D3--D1}) upsample and concatenate
  encoder skip features $s_1$--$s_3$ (blue dashed arrows).
  A dual $3{\times}3$ convolutional head produces the denoised mean $\hat{\mu}_y$
  and aleatoric log-variance $\log\hat{\sigma}^2_a$ simultaneously.
  Epistemic uncertainty $\hat{\sigma}^2_e$ is estimated from the variance of
  $K{=}20$ stochastic forward passes through the VAE Bottleneck (purple branch).
  Coloured grids between blocks represent the feature-map tensor at each
  spatial scale.}
\label{fig:arch:overview}
\end{figure}
\end{landscape}
 from main.tex  (no \documentclass)
\begin{landscape}
\begin{figure}[p]
\centering
\resizebox{0.98\linewidth}{!}{%
\begin{tikzpicture}[
  >=Stealth,
  every node/.style={font=\small},
  % ── block styles ───────────────────────────────────────────────────
  imgblk/.style={inner sep=1pt, outer sep=0pt},
  stemblk/.style={draw=gray!70, fill=cyan!8, rounded corners=3pt,
                   minimum width=2.0cm, minimum height=1.2cm,
                   align=center, line width=0.7pt},
  encblk/.style={draw=blue!55!black, fill=blue!9, rounded corners=3pt,
                  minimum width=2.3cm, minimum height=1.3cm,
                  align=center, line width=0.7pt},
  botblk/.style={draw=teal!70!black, fill=teal!12, rounded corners=3pt,
                  minimum width=2.6cm, minimum height=1.3cm,
                  align=center, line width=1.0pt},
  vaeblk/.style={draw=purple!60!black, fill=purple!7, rounded corners=3pt,
                  minimum width=3.4cm, minimum height=1.6cm,
                  align=center, line width=1.0pt, dashed},
  decblk/.style={draw=orange!65!black, fill=orange!9, rounded corners=3pt,
                  minimum width=2.3cm, minimum height=1.3cm,
                  align=center, line width=0.7pt},
  headblk/.style={draw=teal!55!black, fill=teal!8, rounded corners=3pt,
                   minimum width=1.9cm, minimum height=0.95cm,
                   align=center, line width=0.7pt},
  mcblk/.style={draw=purple!50!black, fill=purple!6, rounded corners=3pt,
                 minimum width=2.4cm, minimum height=0.9cm,
                 align=center, line width=0.8pt, dashed},
  dimtag/.style={font=\fontsize{7}{8}\selectfont, text=gray!65, align=center},
  % ── arrow styles ───────────────────────────────────────────────────
  arrmain/.style={->, line width=1.2pt, gray!55},
  arrskip/.style={->, line width=1.0pt, blue!55, dashed},
  arrbyp/.style={->, line width=0.8pt, gray!40, dashed},
  arrvae/.style={->, line width=0.9pt, purple!55},
]

% ────────────────────────────────────────────────────────────────────
% INPUT THUMBNAIL
% ────────────────────────────────────────────────────────────────────
\node[imgblk] (inp) at (0,8.8)
  {\includegraphics[width=1.6cm,height=1.6cm,keepaspectratio]%
    {figures/arch/arch_input_noisy.png}};
\node[dimtag,below=2pt of inp] {Noisy CXR\\[-1pt]$1{\times}256{\times}256$};

% 1-ch feature grid at input
\begin{scope}[shift={(1.1,8.5)}]
  \fill[gray!25] (0.04,0.04) rectangle (0.84,0.84);
  \fill[gray!35,draw=gray!60,line width=0.4pt] (0,0) rectangle (0.8,0.8);
  \draw[gray!50,very thin,step=0.2] (0,0) grid (0.8,0.8);
\end{scope}

% ────────────────────────────────────────────────────────────────────
% STEM
% ────────────────────────────────────────────────────────────────────
\node[stemblk] (stem) at (2.7,8.8)
  {\textbf{Stem}\\{\fontsize{7}{8}\selectfont Conv $3{\times}3$, pad\,1\\$1\!\to\!64\,\mathrm{ch}$}};

% 64-ch feature map after stem (256²)
\begin{scope}[shift={(3.85,8.5)}]
  \fill[blue!15] (0.04,0.04) rectangle (0.84,0.84);
  \fill[blue!20,draw=blue!45,line width=0.4pt] (0,0) rectangle (0.8,0.8);
  \draw[blue!35,very thin,step=0.2] (0,0) grid (0.8,0.8);
  \node[dimtag,above=0pt] at (0.4,0.85) {$64{\times}256^2$};
\end{scope}

% ────────────────────────────────────────────────────────────────────
% ENCODER
% ────────────────────────────────────────────────────────────────────
\node[encblk] (enc1) at (5.2,8.8)
  {\textbf{E1}\\{\fontsize{7}{8}\selectfont HB${\times}2+$Down}};
\node[dimtag,above=2pt of enc1] {$64{\times}256^2$};

% 128-ch feature map after E1 + downsample (128²)
\begin{scope}[shift={(6.45,8.5)}]
  \fill[blue!18] (0.04,0.04) rectangle (0.69,0.69);
  \fill[blue!24,draw=blue!50,line width=0.4pt] (0,0) rectangle (0.65,0.65);
  \draw[blue!40,very thin,step=0.1625] (0,0) grid (0.65,0.65);
  \node[dimtag,above=0pt] at (0.325,0.7) {$128{\times}128^2$};
\end{scope}

\node[encblk] (enc2) at (7.4,6.3)
  {\textbf{E2}\\{\fontsize{7}{8}\selectfont HB${\times}2+$Down}};
\node[dimtag,above left=1pt and 0pt of enc2] {$128{\times}128^2$};

% 256-ch feature map after E2 + downsample (64²)
\begin{scope}[shift={(8.65,6.05)}]
  \fill[blue!22] (0.04,0.04) rectangle (0.54,0.54);
  \fill[blue!28,draw=blue!55,line width=0.4pt] (0,0) rectangle (0.5,0.5);
  \draw[blue!45,very thin,step=0.125] (0,0) grid (0.5,0.5);
  \node[dimtag,above=0pt] at (0.25,0.55) {$256{\times}64^2$};
\end{scope}

\node[encblk] (enc3) at (9.6,3.8)
  {\textbf{E3}\\{\fontsize{7}{8}\selectfont HB${\times}2+$Down}};
\node[dimtag,above left=1pt and 0pt of enc3] {$256{\times}64^2$};

% 512-ch feature map after E3 + downsample (32²)
\begin{scope}[shift={(10.85,3.6)}]
  \fill[teal!20] (0.04,0.04) rectangle (0.42,0.42);
  \fill[teal!26,draw=teal!55,line width=0.4pt] (0,0) rectangle (0.38,0.38);
  \draw[teal!45,very thin,step=0.095] (0,0) grid (0.38,0.38);
  \node[dimtag,above=0pt] at (0.19,0.43) {$512{\times}32^2$};
\end{scope}

% ────────────────────────────────────────────────────────────────────
% BOTTLENECK
% ────────────────────────────────────────────────────────────────────
\node[botblk] (bot) at (11.9,1.3)
  {\textbf{Bottleneck}\\{\fontsize{7}{8}\selectfont HB${\times}2$}};
\node[dimtag,below=2pt of bot] {$512{\times}32^2$};

% 512-ch feature map after bottleneck
\begin{scope}[shift={(13.3,1.1)}]
  \fill[teal!20] (0.04,0.04) rectangle (0.42,0.42);
  \fill[teal!26,draw=teal!55,line width=0.4pt] (0,0) rectangle (0.38,0.38);
  \draw[teal!45,very thin,step=0.095] (0,0) grid (0.38,0.38);
\end{scope}

% ────────────────────────────────────────────────────────────────────
% VAE BOTTLENECK
% ────────────────────────────────────────────────────────────────────
\node[vaeblk] (vae) at (15.6,1.3)
  {\textbf{VAE Bottleneck}\\
   {\fontsize{7}{8}\selectfont
    $\mathbf{z}_\mu,\,\mathbf{z}_{\log\sigma^2}$\;(Conv $1{\times}1$)\\[2pt]
    $\mathbf{z}{=}\mathbf{z}_\mu{+}\boldsymbol{\varepsilon}{\odot}
     \exp\!\bigl(\tfrac{1}{2}\mathbf{z}_{\log\sigma^2}\bigr)$\\
    $\boldsymbol{\varepsilon}{\sim}\mathcal{N}(\mathbf{0},\mathbf{I})$}};

% ────────────────────────────────────────────────────────────────────
% DECODER
% ────────────────────────────────────────────────────────────────────
\node[decblk] (dec3) at (18.8,3.8)
  {\textbf{D3}\\{\fontsize{7}{8}\selectfont Up$+s_3+$HB${\times}2$}};
\node[dimtag,above right=1pt and 0pt of dec3] {$256{\times}64^2$};

% 256-ch feature map after D3
\begin{scope}[shift={(20.05,3.6)}]
  \fill[orange!20] (0.04,0.04) rectangle (0.54,0.54);
  \fill[orange!26,draw=orange!55,line width=0.4pt] (0,0) rectangle (0.5,0.5);
  \draw[orange!45,very thin,step=0.125] (0,0) grid (0.5,0.5);
  \node[dimtag,above=0pt] at (0.25,0.55) {$256{\times}64^2$};
\end{scope}

\node[decblk] (dec2) at (21.0,6.3)
  {\textbf{D2}\\{\fontsize{7}{8}\selectfont Up$+s_2+$HB${\times}2$}};
\node[dimtag,above right=1pt and 0pt of dec2] {$128{\times}128^2$};

% 128-ch feature map after D2
\begin{scope}[shift={(22.25,6.05)}]
  \fill[orange!18] (0.04,0.04) rectangle (0.69,0.69);
  \fill[orange!23,draw=orange!50,line width=0.4pt] (0,0) rectangle (0.65,0.65);
  \draw[orange!40,very thin,step=0.1625] (0,0) grid (0.65,0.65);
  \node[dimtag,above=0pt] at (0.325,0.7) {$128{\times}128^2$};
\end{scope}

\node[decblk] (dec1) at (23.2,8.8)
  {\textbf{D1}\\{\fontsize{7}{8}\selectfont Up$+s_1+$HB${\times}2$}};
\node[dimtag,above=2pt of dec1] {$64{\times}256^2$};

% ────────────────────────────────────────────────────────────────────
% DUAL HEAD + OUTPUT THUMBNAILS
% ────────────────────────────────────────────────────────────────────
\node[headblk] (hmu) at (25.2,9.5)
  {$\hat{\mu}_y$\\{\fontsize{7}{8}\selectfont Conv $3{\times}3$}};
\node[headblk,draw=orange!60!black,fill=orange!8] (hsig) at (25.2,8.1)
  {$\log\hat{\sigma}^2_a$\\{\fontsize{7}{8}\selectfont Conv $3{\times}3$}};
\node[dimtag,above=2pt of hmu] {Dual Head};

\node[imgblk] (outmu) at (27.0,9.5)
  {\includegraphics[width=1.4cm,height=1.4cm,keepaspectratio]%
    {figures/arch/arch_output_mu.png}};
\node[dimtag,below=1pt of outmu] {$\hat{\mu}_y$};

\node[imgblk] (outsig) at (27.0,8.1)
  {\includegraphics[width=1.4cm,height=1.4cm,keepaspectratio]%
    {figures/arch/arch_output_aleat.png}};
\node[dimtag,below=1pt of outsig] {$\hat{\sigma}^2_a$};

% ────────────────────────────────────────────────────────────────────
% K=20 EPISTEMIC BRANCH
% ────────────────────────────────────────────────────────────────────
\node[mcblk] (mc) at (15.6,-0.3)
  {{\fontsize{7}{8}\selectfont $K{=}20$ stochastic passes}};
\node[imgblk] (outepi) at (18.8,-0.3)
  {\includegraphics[width=1.4cm,height=1.4cm,keepaspectratio]%
    {figures/arch/arch_output_epist.png}};
\node[dimtag,below=1pt of outepi] {$\hat{\sigma}^2_e$\,(epistemic)};

% ────────────────────────────────────────────────────────────────────
% ARROWS – main flow
% ────────────────────────────────────────────────────────────────────
\draw[arrmain] (inp.east)          -- (stem.west);
\draw[arrmain] (stem.east)         -- (enc1.west);
\draw[arrmain] (enc1.south east)   -- (enc2.north west);
\draw[arrmain] (enc2.south east)   -- (enc3.north west);
\draw[arrmain] (enc3.south east)   -- (bot.north west);
\draw[arrmain] (bot.east)          -- (vae.west);
\draw[arrmain] (vae.east)          -| (dec3.south);
\draw[arrmain] (dec3.north east)   -- (dec2.south west);
\draw[arrmain] (dec2.north east)   -- (dec1.south west);
\draw[arrmain] (dec1.east) -- ++(0.4,0) |- (hmu.west);
\draw[arrmain] (dec1.east) -- ++(0.4,0) |- (hsig.west);
\draw[arrmain] (hmu.east)          -- (outmu.west);
\draw[arrmain] (hsig.east)         -- (outsig.west);

% Arms A-D deterministic bypass (skips VAE)
\draw[arrbyp] (bot.north) -- ++(0,1.5) -|
  node[above,pos=0.3,font=\fontsize{6}{7}\selectfont,text=gray!55]
    {Arms\,A--D\;(no\,VAE)}
  (dec3.west);

% VAE → epistemic MC branch
\draw[arrvae] (vae.south) -- (mc.north);
\draw[arrvae] (mc.east)   -- (outepi.west);
\node[dimtag,right=2pt of outepi] {Var of $K$ predictions};

% ────────────────────────────────────────────────────────────────────
% SKIP CONNECTIONS (dashed blue, routed above encoder blocks)
% enc1.x=5.2, dec1.x=23.2  → dx=18.0, route at y=10.05
% enc2.x=7.4, dec2.x=21.0  → dx=13.6, route at y=7.4
% enc3.x=9.6, dec3.x=18.8  → dx=9.2,  route at y=4.85
% ────────────────────────────────────────────────────────────────────
\draw[arrskip]
  (enc1.north) -- ++(0,0.6) --
  node[above,font=\fontsize{7}{8}\selectfont,black!70]{$s_1$}
  ++(18.0,0) -- (dec1.north);

\draw[arrskip]
  (enc2.north) -- ++(0,0.45) --
  node[above,font=\fontsize{7}{8}\selectfont,black!70]{$s_2$}
  ++(13.6,0) -- (dec2.north);

\draw[arrskip]
  (enc3.north) -- ++(0,0.4) --
  node[above,font=\fontsize{7}{8}\selectfont,black!70]{$s_3$}
  ++(9.2,0) -- (dec3.north);

\end{tikzpicture}%
}
\caption[PP-VAE-Hformer overall architecture]{%
  \textbf{PP-VAE-Hformer encoder--decoder architecture.}
  A noisy paediatric CXR ($1{\times}256{\times}256$) enters a stem convolution
  and three encoder stages (\textbf{E1--E3}), each comprising two HybridBlocks
  (HB${\times}2$) followed by strided downsampling; channel depth expands
  $64\!\to\!128\!\to\!256\,\text{ch}$ while spatial resolution halves.
  The Bottleneck applies two further HybridBlocks at $512{\times}32^2$.
  The probabilistic VAE Bottleneck (dashed border) maps features to a spatial
  latent via the reparameterisation trick; Arms\,A--D bypass this stage
  (grey dashed arrow).
  Three symmetric decoder stages (\textbf{D3--D1}) upsample and concatenate
  encoder skip features $s_1$--$s_3$ (blue dashed arrows).
  A dual $3{\times}3$ convolutional head produces the denoised mean $\hat{\mu}_y$
  and aleatoric log-variance $\log\hat{\sigma}^2_a$ simultaneously.
  Epistemic uncertainty $\hat{\sigma}^2_e$ is estimated from the variance of
  $K{=}20$ stochastic forward passes through the VAE Bottleneck (purple branch).
  Coloured grids between blocks represent the feature-map tensor at each
  spatial scale.}
\label{fig:arch:overview}
\end{figure}
\end{landscape}
 from main.tex  (no \documentclass)
\begin{landscape}
\begin{figure}[p]
\centering
\resizebox{0.98\linewidth}{!}{%
\begin{tikzpicture}[
  >=Stealth,
  every node/.style={font=\small},
  % ── block styles ───────────────────────────────────────────────────
  imgblk/.style={inner sep=1pt, outer sep=0pt},
  stemblk/.style={draw=gray!70, fill=cyan!8, rounded corners=3pt,
                   minimum width=2.0cm, minimum height=1.2cm,
                   align=center, line width=0.7pt},
  encblk/.style={draw=blue!55!black, fill=blue!9, rounded corners=3pt,
                  minimum width=2.3cm, minimum height=1.3cm,
                  align=center, line width=0.7pt},
  botblk/.style={draw=teal!70!black, fill=teal!12, rounded corners=3pt,
                  minimum width=2.6cm, minimum height=1.3cm,
                  align=center, line width=1.0pt},
  vaeblk/.style={draw=purple!60!black, fill=purple!7, rounded corners=3pt,
                  minimum width=3.4cm, minimum height=1.6cm,
                  align=center, line width=1.0pt, dashed},
  decblk/.style={draw=orange!65!black, fill=orange!9, rounded corners=3pt,
                  minimum width=2.3cm, minimum height=1.3cm,
                  align=center, line width=0.7pt},
  headblk/.style={draw=teal!55!black, fill=teal!8, rounded corners=3pt,
                   minimum width=1.9cm, minimum height=0.95cm,
                   align=center, line width=0.7pt},
  mcblk/.style={draw=purple!50!black, fill=purple!6, rounded corners=3pt,
                 minimum width=2.4cm, minimum height=0.9cm,
                 align=center, line width=0.8pt, dashed},
  dimtag/.style={font=\fontsize{7}{8}\selectfont, text=gray!65, align=center},
  % ── arrow styles ───────────────────────────────────────────────────
  arrmain/.style={->, line width=1.2pt, gray!55},
  arrskip/.style={->, line width=1.0pt, blue!55, dashed},
  arrbyp/.style={->, line width=0.8pt, gray!40, dashed},
  arrvae/.style={->, line width=0.9pt, purple!55},
]

% ────────────────────────────────────────────────────────────────────
% INPUT THUMBNAIL
% ────────────────────────────────────────────────────────────────────
\node[imgblk] (inp) at (0,8.8)
  {\includegraphics[width=1.6cm,height=1.6cm,keepaspectratio]%
    {figures/arch/arch_input_noisy.png}};
\node[dimtag,below=2pt of inp] {Noisy CXR\\[-1pt]$1{\times}256{\times}256$};

% 1-ch feature grid at input
\begin{scope}[shift={(1.1,8.5)}]
  \fill[gray!25] (0.04,0.04) rectangle (0.84,0.84);
  \fill[gray!35,draw=gray!60,line width=0.4pt] (0,0) rectangle (0.8,0.8);
  \draw[gray!50,very thin,step=0.2] (0,0) grid (0.8,0.8);
\end{scope}

% ────────────────────────────────────────────────────────────────────
% STEM
% ────────────────────────────────────────────────────────────────────
\node[stemblk] (stem) at (2.7,8.8)
  {\textbf{Stem}\\{\fontsize{7}{8}\selectfont Conv $3{\times}3$, pad\,1\\$1\!\to\!64\,\mathrm{ch}$}};

% 64-ch feature map after stem (256²)
\begin{scope}[shift={(3.85,8.5)}]
  \fill[blue!15] (0.04,0.04) rectangle (0.84,0.84);
  \fill[blue!20,draw=blue!45,line width=0.4pt] (0,0) rectangle (0.8,0.8);
  \draw[blue!35,very thin,step=0.2] (0,0) grid (0.8,0.8);
  \node[dimtag,above=0pt] at (0.4,0.85) {$64{\times}256^2$};
\end{scope}

% ────────────────────────────────────────────────────────────────────
% ENCODER
% ────────────────────────────────────────────────────────────────────
\node[encblk] (enc1) at (5.2,8.8)
  {\textbf{E1}\\{\fontsize{7}{8}\selectfont HB${\times}2+$Down}};
\node[dimtag,above=2pt of enc1] {$64{\times}256^2$};

% 128-ch feature map after E1 + downsample (128²)
\begin{scope}[shift={(6.45,8.5)}]
  \fill[blue!18] (0.04,0.04) rectangle (0.69,0.69);
  \fill[blue!24,draw=blue!50,line width=0.4pt] (0,0) rectangle (0.65,0.65);
  \draw[blue!40,very thin,step=0.1625] (0,0) grid (0.65,0.65);
  \node[dimtag,above=0pt] at (0.325,0.7) {$128{\times}128^2$};
\end{scope}

\node[encblk] (enc2) at (7.4,6.3)
  {\textbf{E2}\\{\fontsize{7}{8}\selectfont HB${\times}2+$Down}};
\node[dimtag,above left=1pt and 0pt of enc2] {$128{\times}128^2$};

% 256-ch feature map after E2 + downsample (64²)
\begin{scope}[shift={(8.65,6.05)}]
  \fill[blue!22] (0.04,0.04) rectangle (0.54,0.54);
  \fill[blue!28,draw=blue!55,line width=0.4pt] (0,0) rectangle (0.5,0.5);
  \draw[blue!45,very thin,step=0.125] (0,0) grid (0.5,0.5);
  \node[dimtag,above=0pt] at (0.25,0.55) {$256{\times}64^2$};
\end{scope}

\node[encblk] (enc3) at (9.6,3.8)
  {\textbf{E3}\\{\fontsize{7}{8}\selectfont HB${\times}2+$Down}};
\node[dimtag,above left=1pt and 0pt of enc3] {$256{\times}64^2$};

% 512-ch feature map after E3 + downsample (32²)
\begin{scope}[shift={(10.85,3.6)}]
  \fill[teal!20] (0.04,0.04) rectangle (0.42,0.42);
  \fill[teal!26,draw=teal!55,line width=0.4pt] (0,0) rectangle (0.38,0.38);
  \draw[teal!45,very thin,step=0.095] (0,0) grid (0.38,0.38);
  \node[dimtag,above=0pt] at (0.19,0.43) {$512{\times}32^2$};
\end{scope}

% ────────────────────────────────────────────────────────────────────
% BOTTLENECK
% ────────────────────────────────────────────────────────────────────
\node[botblk] (bot) at (11.9,1.3)
  {\textbf{Bottleneck}\\{\fontsize{7}{8}\selectfont HB${\times}2$}};
\node[dimtag,below=2pt of bot] {$512{\times}32^2$};

% 512-ch feature map after bottleneck
\begin{scope}[shift={(13.3,1.1)}]
  \fill[teal!20] (0.04,0.04) rectangle (0.42,0.42);
  \fill[teal!26,draw=teal!55,line width=0.4pt] (0,0) rectangle (0.38,0.38);
  \draw[teal!45,very thin,step=0.095] (0,0) grid (0.38,0.38);
\end{scope}

% ────────────────────────────────────────────────────────────────────
% VAE BOTTLENECK
% ────────────────────────────────────────────────────────────────────
\node[vaeblk] (vae) at (15.6,1.3)
  {\textbf{VAE Bottleneck}\\
   {\fontsize{7}{8}\selectfont
    $\mathbf{z}_\mu,\,\mathbf{z}_{\log\sigma^2}$\;(Conv $1{\times}1$)\\[2pt]
    $\mathbf{z}{=}\mathbf{z}_\mu{+}\boldsymbol{\varepsilon}{\odot}
     \exp\!\bigl(\tfrac{1}{2}\mathbf{z}_{\log\sigma^2}\bigr)$\\
    $\boldsymbol{\varepsilon}{\sim}\mathcal{N}(\mathbf{0},\mathbf{I})$}};

% ────────────────────────────────────────────────────────────────────
% DECODER
% ────────────────────────────────────────────────────────────────────
\node[decblk] (dec3) at (18.8,3.8)
  {\textbf{D3}\\{\fontsize{7}{8}\selectfont Up$+s_3+$HB${\times}2$}};
\node[dimtag,above right=1pt and 0pt of dec3] {$256{\times}64^2$};

% 256-ch feature map after D3
\begin{scope}[shift={(20.05,3.6)}]
  \fill[orange!20] (0.04,0.04) rectangle (0.54,0.54);
  \fill[orange!26,draw=orange!55,line width=0.4pt] (0,0) rectangle (0.5,0.5);
  \draw[orange!45,very thin,step=0.125] (0,0) grid (0.5,0.5);
  \node[dimtag,above=0pt] at (0.25,0.55) {$256{\times}64^2$};
\end{scope}

\node[decblk] (dec2) at (21.0,6.3)
  {\textbf{D2}\\{\fontsize{7}{8}\selectfont Up$+s_2+$HB${\times}2$}};
\node[dimtag,above right=1pt and 0pt of dec2] {$128{\times}128^2$};

% 128-ch feature map after D2
\begin{scope}[shift={(22.25,6.05)}]
  \fill[orange!18] (0.04,0.04) rectangle (0.69,0.69);
  \fill[orange!23,draw=orange!50,line width=0.4pt] (0,0) rectangle (0.65,0.65);
  \draw[orange!40,very thin,step=0.1625] (0,0) grid (0.65,0.65);
  \node[dimtag,above=0pt] at (0.325,0.7) {$128{\times}128^2$};
\end{scope}

\node[decblk] (dec1) at (23.2,8.8)
  {\textbf{D1}\\{\fontsize{7}{8}\selectfont Up$+s_1+$HB${\times}2$}};
\node[dimtag,above=2pt of dec1] {$64{\times}256^2$};

% ────────────────────────────────────────────────────────────────────
% DUAL HEAD + OUTPUT THUMBNAILS
% ────────────────────────────────────────────────────────────────────
\node[headblk] (hmu) at (25.2,9.5)
  {$\hat{\mu}_y$\\{\fontsize{7}{8}\selectfont Conv $3{\times}3$}};
\node[headblk,draw=orange!60!black,fill=orange!8] (hsig) at (25.2,8.1)
  {$\log\hat{\sigma}^2_a$\\{\fontsize{7}{8}\selectfont Conv $3{\times}3$}};
\node[dimtag,above=2pt of hmu] {Dual Head};

\node[imgblk] (outmu) at (27.0,9.5)
  {\includegraphics[width=1.4cm,height=1.4cm,keepaspectratio]%
    {figures/arch/arch_output_mu.png}};
\node[dimtag,below=1pt of outmu] {$\hat{\mu}_y$};

\node[imgblk] (outsig) at (27.0,8.1)
  {\includegraphics[width=1.4cm,height=1.4cm,keepaspectratio]%
    {figures/arch/arch_output_aleat.png}};
\node[dimtag,below=1pt of outsig] {$\hat{\sigma}^2_a$};

% ────────────────────────────────────────────────────────────────────
% K=20 EPISTEMIC BRANCH
% ────────────────────────────────────────────────────────────────────
\node[mcblk] (mc) at (15.6,-0.3)
  {{\fontsize{7}{8}\selectfont $K{=}20$ stochastic passes}};
\node[imgblk] (outepi) at (18.8,-0.3)
  {\includegraphics[width=1.4cm,height=1.4cm,keepaspectratio]%
    {figures/arch/arch_output_epist.png}};
\node[dimtag,below=1pt of outepi] {$\hat{\sigma}^2_e$\,(epistemic)};

% ────────────────────────────────────────────────────────────────────
% ARROWS – main flow
% ────────────────────────────────────────────────────────────────────
\draw[arrmain] (inp.east)          -- (stem.west);
\draw[arrmain] (stem.east)         -- (enc1.west);
\draw[arrmain] (enc1.south east)   -- (enc2.north west);
\draw[arrmain] (enc2.south east)   -- (enc3.north west);
\draw[arrmain] (enc3.south east)   -- (bot.north west);
\draw[arrmain] (bot.east)          -- (vae.west);
\draw[arrmain] (vae.east)          -| (dec3.south);
\draw[arrmain] (dec3.north east)   -- (dec2.south west);
\draw[arrmain] (dec2.north east)   -- (dec1.south west);
\draw[arrmain] (dec1.east) -- ++(0.4,0) |- (hmu.west);
\draw[arrmain] (dec1.east) -- ++(0.4,0) |- (hsig.west);
\draw[arrmain] (hmu.east)          -- (outmu.west);
\draw[arrmain] (hsig.east)         -- (outsig.west);

% Arms A-D deterministic bypass (skips VAE)
\draw[arrbyp] (bot.north) -- ++(0,1.5) -|
  node[above,pos=0.3,font=\fontsize{6}{7}\selectfont,text=gray!55]
    {Arms\,A--D\;(no\,VAE)}
  (dec3.west);

% VAE → epistemic MC branch
\draw[arrvae] (vae.south) -- (mc.north);
\draw[arrvae] (mc.east)   -- (outepi.west);
\node[dimtag,right=2pt of outepi] {Var of $K$ predictions};

% ────────────────────────────────────────────────────────────────────
% SKIP CONNECTIONS (dashed blue, routed above encoder blocks)
% enc1.x=5.2, dec1.x=23.2  → dx=18.0, route at y=10.05
% enc2.x=7.4, dec2.x=21.0  → dx=13.6, route at y=7.4
% enc3.x=9.6, dec3.x=18.8  → dx=9.2,  route at y=4.85
% ────────────────────────────────────────────────────────────────────
\draw[arrskip]
  (enc1.north) -- ++(0,0.6) --
  node[above,font=\fontsize{7}{8}\selectfont,black!70]{$s_1$}
  ++(18.0,0) -- (dec1.north);

\draw[arrskip]
  (enc2.north) -- ++(0,0.45) --
  node[above,font=\fontsize{7}{8}\selectfont,black!70]{$s_2$}
  ++(13.6,0) -- (dec2.north);

\draw[arrskip]
  (enc3.north) -- ++(0,0.4) --
  node[above,font=\fontsize{7}{8}\selectfont,black!70]{$s_3$}
  ++(9.2,0) -- (dec3.north);

\end{tikzpicture}%
}
\caption[PP-VAE-Hformer overall architecture]{%
  \textbf{PP-VAE-Hformer encoder--decoder architecture.}
  A noisy paediatric CXR ($1{\times}256{\times}256$) enters a stem convolution
  and three encoder stages (\textbf{E1--E3}), each comprising two HybridBlocks
  (HB${\times}2$) followed by strided downsampling; channel depth expands
  $64\!\to\!128\!\to\!256\,\text{ch}$ while spatial resolution halves.
  The Bottleneck applies two further HybridBlocks at $512{\times}32^2$.
  The probabilistic VAE Bottleneck (dashed border) maps features to a spatial
  latent via the reparameterisation trick; Arms\,A--D bypass this stage
  (grey dashed arrow).
  Three symmetric decoder stages (\textbf{D3--D1}) upsample and concatenate
  encoder skip features $s_1$--$s_3$ (blue dashed arrows).
  A dual $3{\times}3$ convolutional head produces the denoised mean $\hat{\mu}_y$
  and aleatoric log-variance $\log\hat{\sigma}^2_a$ simultaneously.
  Epistemic uncertainty $\hat{\sigma}^2_e$ is estimated from the variance of
  $K{=}20$ stochastic forward passes through the VAE Bottleneck (purple branch).
  Coloured grids between blocks represent the feature-map tensor at each
  spatial scale.}
\label{fig:arch:overview}
\end{figure}
\end{landscape}
 from main.tex  (no \documentclass)
\begin{landscape}
\begin{figure}[p]
\centering
\resizebox{0.98\linewidth}{!}{%
\begin{tikzpicture}[
  >=Stealth,
  every node/.style={font=\small},
  % ── block styles ───────────────────────────────────────────────────
  imgblk/.style={inner sep=1pt, outer sep=0pt},
  stemblk/.style={draw=gray!70, fill=cyan!8, rounded corners=3pt,
                   minimum width=2.0cm, minimum height=1.2cm,
                   align=center, line width=0.7pt},
  encblk/.style={draw=blue!55!black, fill=blue!9, rounded corners=3pt,
                  minimum width=2.3cm, minimum height=1.3cm,
                  align=center, line width=0.7pt},
  botblk/.style={draw=teal!70!black, fill=teal!12, rounded corners=3pt,
                  minimum width=2.6cm, minimum height=1.3cm,
                  align=center, line width=1.0pt},
  vaeblk/.style={draw=purple!60!black, fill=purple!7, rounded corners=3pt,
                  minimum width=3.4cm, minimum height=1.6cm,
                  align=center, line width=1.0pt, dashed},
  decblk/.style={draw=orange!65!black, fill=orange!9, rounded corners=3pt,
                  minimum width=2.3cm, minimum height=1.3cm,
                  align=center, line width=0.7pt},
  headblk/.style={draw=teal!55!black, fill=teal!8, rounded corners=3pt,
                   minimum width=1.9cm, minimum height=0.95cm,
                   align=center, line width=0.7pt},
  mcblk/.style={draw=purple!50!black, fill=purple!6, rounded corners=3pt,
                 minimum width=2.4cm, minimum height=0.9cm,
                 align=center, line width=0.8pt, dashed},
  dimtag/.style={font=\fontsize{7}{8}\selectfont, text=gray!65, align=center},
  % ── arrow styles ───────────────────────────────────────────────────
  arrmain/.style={->, line width=1.2pt, gray!55},
  arrskip/.style={->, line width=1.0pt, blue!55, dashed},
  arrbyp/.style={->, line width=0.8pt, gray!40, dashed},
  arrvae/.style={->, line width=0.9pt, purple!55},
]

% ────────────────────────────────────────────────────────────────────
% INPUT THUMBNAIL
% ────────────────────────────────────────────────────────────────────
\node[imgblk] (inp) at (0,8.8)
  {\includegraphics[width=1.6cm,height=1.6cm,keepaspectratio]%
    {figures/arch/arch_input_noisy.png}};
\node[dimtag,below=2pt of inp] {Noisy CXR\\[-1pt]$1{\times}256{\times}256$};

% 1-ch feature grid at input
\begin{scope}[shift={(1.1,8.5)}]
  \fill[gray!25] (0.04,0.04) rectangle (0.84,0.84);
  \fill[gray!35,draw=gray!60,line width=0.4pt] (0,0) rectangle (0.8,0.8);
  \draw[gray!50,very thin,step=0.2] (0,0) grid (0.8,0.8);
\end{scope}

% ────────────────────────────────────────────────────────────────────
% STEM
% ────────────────────────────────────────────────────────────────────
\node[stemblk] (stem) at (2.7,8.8)
  {\textbf{Stem}\\{\fontsize{7}{8}\selectfont Conv $3{\times}3$, pad\,1\\$1\!\to\!64\,\mathrm{ch}$}};

% 64-ch feature map after stem (256²)
\begin{scope}[shift={(3.85,8.5)}]
  \fill[blue!15] (0.04,0.04) rectangle (0.84,0.84);
  \fill[blue!20,draw=blue!45,line width=0.4pt] (0,0) rectangle (0.8,0.8);
  \draw[blue!35,very thin,step=0.2] (0,0) grid (0.8,0.8);
  \node[dimtag,above=0pt] at (0.4,0.85) {$64{\times}256^2$};
\end{scope}

% ────────────────────────────────────────────────────────────────────
% ENCODER
% ────────────────────────────────────────────────────────────────────
\node[encblk] (enc1) at (5.2,8.8)
  {\textbf{E1}\\{\fontsize{7}{8}\selectfont HB${\times}2+$Down}};
\node[dimtag,above=2pt of enc1] {$64{\times}256^2$};

% 128-ch feature map after E1 + downsample (128²)
\begin{scope}[shift={(6.45,8.5)}]
  \fill[blue!18] (0.04,0.04) rectangle (0.69,0.69);
  \fill[blue!24,draw=blue!50,line width=0.4pt] (0,0) rectangle (0.65,0.65);
  \draw[blue!40,very thin,step=0.1625] (0,0) grid (0.65,0.65);
  \node[dimtag,above=0pt] at (0.325,0.7) {$128{\times}128^2$};
\end{scope}

\node[encblk] (enc2) at (7.4,6.3)
  {\textbf{E2}\\{\fontsize{7}{8}\selectfont HB${\times}2+$Down}};
\node[dimtag,above left=1pt and 0pt of enc2] {$128{\times}128^2$};

% 256-ch feature map after E2 + downsample (64²)
\begin{scope}[shift={(8.65,6.05)}]
  \fill[blue!22] (0.04,0.04) rectangle (0.54,0.54);
  \fill[blue!28,draw=blue!55,line width=0.4pt] (0,0) rectangle (0.5,0.5);
  \draw[blue!45,very thin,step=0.125] (0,0) grid (0.5,0.5);
  \node[dimtag,above=0pt] at (0.25,0.55) {$256{\times}64^2$};
\end{scope}

\node[encblk] (enc3) at (9.6,3.8)
  {\textbf{E3}\\{\fontsize{7}{8}\selectfont HB${\times}2+$Down}};
\node[dimtag,above left=1pt and 0pt of enc3] {$256{\times}64^2$};

% 512-ch feature map after E3 + downsample (32²)
\begin{scope}[shift={(10.85,3.6)}]
  \fill[teal!20] (0.04,0.04) rectangle (0.42,0.42);
  \fill[teal!26,draw=teal!55,line width=0.4pt] (0,0) rectangle (0.38,0.38);
  \draw[teal!45,very thin,step=0.095] (0,0) grid (0.38,0.38);
  \node[dimtag,above=0pt] at (0.19,0.43) {$512{\times}32^2$};
\end{scope}

% ────────────────────────────────────────────────────────────────────
% BOTTLENECK
% ────────────────────────────────────────────────────────────────────
\node[botblk] (bot) at (11.9,1.3)
  {\textbf{Bottleneck}\\{\fontsize{7}{8}\selectfont HB${\times}2$}};
\node[dimtag,below=2pt of bot] {$512{\times}32^2$};

% 512-ch feature map after bottleneck
\begin{scope}[shift={(13.3,1.1)}]
  \fill[teal!20] (0.04,0.04) rectangle (0.42,0.42);
  \fill[teal!26,draw=teal!55,line width=0.4pt] (0,0) rectangle (0.38,0.38);
  \draw[teal!45,very thin,step=0.095] (0,0) grid (0.38,0.38);
\end{scope}

% ────────────────────────────────────────────────────────────────────
% VAE BOTTLENECK
% ────────────────────────────────────────────────────────────────────
\node[vaeblk] (vae) at (15.6,1.3)
  {\textbf{VAE Bottleneck}\\
   {\fontsize{7}{8}\selectfont
    $\mathbf{z}_\mu,\,\mathbf{z}_{\log\sigma^2}$\;(Conv $1{\times}1$)\\[2pt]
    $\mathbf{z}{=}\mathbf{z}_\mu{+}\boldsymbol{\varepsilon}{\odot}
     \exp\!\bigl(\tfrac{1}{2}\mathbf{z}_{\log\sigma^2}\bigr)$\\
    $\boldsymbol{\varepsilon}{\sim}\mathcal{N}(\mathbf{0},\mathbf{I})$}};

% ────────────────────────────────────────────────────────────────────
% DECODER
% ────────────────────────────────────────────────────────────────────
\node[decblk] (dec3) at (18.8,3.8)
  {\textbf{D3}\\{\fontsize{7}{8}\selectfont Up$+s_3+$HB${\times}2$}};
\node[dimtag,above right=1pt and 0pt of dec3] {$256{\times}64^2$};

% 256-ch feature map after D3
\begin{scope}[shift={(20.05,3.6)}]
  \fill[orange!20] (0.04,0.04) rectangle (0.54,0.54);
  \fill[orange!26,draw=orange!55,line width=0.4pt] (0,0) rectangle (0.5,0.5);
  \draw[orange!45,very thin,step=0.125] (0,0) grid (0.5,0.5);
  \node[dimtag,above=0pt] at (0.25,0.55) {$256{\times}64^2$};
\end{scope}

\node[decblk] (dec2) at (21.0,6.3)
  {\textbf{D2}\\{\fontsize{7}{8}\selectfont Up$+s_2+$HB${\times}2$}};
\node[dimtag,above right=1pt and 0pt of dec2] {$128{\times}128^2$};

% 128-ch feature map after D2
\begin{scope}[shift={(22.25,6.05)}]
  \fill[orange!18] (0.04,0.04) rectangle (0.69,0.69);
  \fill[orange!23,draw=orange!50,line width=0.4pt] (0,0) rectangle (0.65,0.65);
  \draw[orange!40,very thin,step=0.1625] (0,0) grid (0.65,0.65);
  \node[dimtag,above=0pt] at (0.325,0.7) {$128{\times}128^2$};
\end{scope}

\node[decblk] (dec1) at (23.2,8.8)
  {\textbf{D1}\\{\fontsize{7}{8}\selectfont Up$+s_1+$HB${\times}2$}};
\node[dimtag,above=2pt of dec1] {$64{\times}256^2$};

% ────────────────────────────────────────────────────────────────────
% DUAL HEAD + OUTPUT THUMBNAILS
% ────────────────────────────────────────────────────────────────────
\node[headblk] (hmu) at (25.2,9.5)
  {$\hat{\mu}_y$\\{\fontsize{7}{8}\selectfont Conv $3{\times}3$}};
\node[headblk,draw=orange!60!black,fill=orange!8] (hsig) at (25.2,8.1)
  {$\log\hat{\sigma}^2_a$\\{\fontsize{7}{8}\selectfont Conv $3{\times}3$}};
\node[dimtag,above=2pt of hmu] {Dual Head};

\node[imgblk] (outmu) at (27.0,9.5)
  {\includegraphics[width=1.4cm,height=1.4cm,keepaspectratio]%
    {figures/arch/arch_output_mu.png}};
\node[dimtag,below=1pt of outmu] {$\hat{\mu}_y$};

\node[imgblk] (outsig) at (27.0,8.1)
  {\includegraphics[width=1.4cm,height=1.4cm,keepaspectratio]%
    {figures/arch/arch_output_aleat.png}};
\node[dimtag,below=1pt of outsig] {$\hat{\sigma}^2_a$};

% ────────────────────────────────────────────────────────────────────
% K=20 EPISTEMIC BRANCH
% ────────────────────────────────────────────────────────────────────
\node[mcblk] (mc) at (15.6,-0.3)
  {{\fontsize{7}{8}\selectfont $K{=}20$ stochastic passes}};
\node[imgblk] (outepi) at (18.8,-0.3)
  {\includegraphics[width=1.4cm,height=1.4cm,keepaspectratio]%
    {figures/arch/arch_output_epist.png}};
\node[dimtag,below=1pt of outepi] {$\hat{\sigma}^2_e$\,(epistemic)};

% ────────────────────────────────────────────────────────────────────
% ARROWS – main flow
% ────────────────────────────────────────────────────────────────────
\draw[arrmain] (inp.east)          -- (stem.west);
\draw[arrmain] (stem.east)         -- (enc1.west);
\draw[arrmain] (enc1.south east)   -- (enc2.north west);
\draw[arrmain] (enc2.south east)   -- (enc3.north west);
\draw[arrmain] (enc3.south east)   -- (bot.north west);
\draw[arrmain] (bot.east)          -- (vae.west);
\draw[arrmain] (vae.east)          -| (dec3.south);
\draw[arrmain] (dec3.north east)   -- (dec2.south west);
\draw[arrmain] (dec2.north east)   -- (dec1.south west);
\draw[arrmain] (dec1.east) -- ++(0.4,0) |- (hmu.west);
\draw[arrmain] (dec1.east) -- ++(0.4,0) |- (hsig.west);
\draw[arrmain] (hmu.east)          -- (outmu.west);
\draw[arrmain] (hsig.east)         -- (outsig.west);

% Arms A-D deterministic bypass (skips VAE)
\draw[arrbyp] (bot.north) -- ++(0,1.5) -|
  node[above,pos=0.3,font=\fontsize{6}{7}\selectfont,text=gray!55]
    {Arms\,A--D\;(no\,VAE)}
  (dec3.west);

% VAE → epistemic MC branch
\draw[arrvae] (vae.south) -- (mc.north);
\draw[arrvae] (mc.east)   -- (outepi.west);
\node[dimtag,right=2pt of outepi] {Var of $K$ predictions};

% ────────────────────────────────────────────────────────────────────
% SKIP CONNECTIONS (dashed blue, routed above encoder blocks)
% enc1.x=5.2, dec1.x=23.2  → dx=18.0, route at y=10.05
% enc2.x=7.4, dec2.x=21.0  → dx=13.6, route at y=7.4
% enc3.x=9.6, dec3.x=18.8  → dx=9.2,  route at y=4.85
% ────────────────────────────────────────────────────────────────────
\draw[arrskip]
  (enc1.north) -- ++(0,0.6) --
  node[above,font=\fontsize{7}{8}\selectfont,black!70]{$s_1$}
  ++(18.0,0) -- (dec1.north);

\draw[arrskip]
  (enc2.north) -- ++(0,0.45) --
  node[above,font=\fontsize{7}{8}\selectfont,black!70]{$s_2$}
  ++(13.6,0) -- (dec2.north);

\draw[arrskip]
  (enc3.north) -- ++(0,0.4) --
  node[above,font=\fontsize{7}{8}\selectfont,black!70]{$s_3$}
  ++(9.2,0) -- (dec3.north);

\end{tikzpicture}%
}
\caption[PP-VAE-Hformer overall architecture]{%
  \textbf{PP-VAE-Hformer encoder--decoder architecture.}
  A noisy paediatric CXR ($1{\times}256{\times}256$) enters a stem convolution
  and three encoder stages (\textbf{E1--E3}), each comprising two HybridBlocks
  (HB${\times}2$) followed by strided downsampling; channel depth expands
  $64\!\to\!128\!\to\!256\,\text{ch}$ while spatial resolution halves.
  The Bottleneck applies two further HybridBlocks at $512{\times}32^2$.
  The probabilistic VAE Bottleneck (dashed border) maps features to a spatial
  latent via the reparameterisation trick; Arms\,A--D bypass this stage
  (grey dashed arrow).
  Three symmetric decoder stages (\textbf{D3--D1}) upsample and concatenate
  encoder skip features $s_1$--$s_3$ (blue dashed arrows).
  A dual $3{\times}3$ convolutional head produces the denoised mean $\hat{\mu}_y$
  and aleatoric log-variance $\log\hat{\sigma}^2_a$ simultaneously.
  Epistemic uncertainty $\hat{\sigma}^2_e$ is estimated from the variance of
  $K{=}20$ stochastic forward passes through the VAE Bottleneck (purple branch).
  Coloured grids between blocks represent the feature-map tensor at each
  spatial scale.}
\label{fig:arch:overview}
\end{figure}
\end{landscape}
